\documentclass[11pt,a4paper]{article}

\usepackage[UTF8]{ctex}
\usepackage{amsmath,amssymb}
\usepackage{geometry}
\usepackage{enumitem}
\usepackage{booktabs}
\usepackage{xcolor}
\usepackage{hyperref}
\usepackage{tikz}
\usetikzlibrary{trees,positioning,arrows}

\geometry{margin=2.2cm}

\definecolor{ansblue}{RGB}{0,70,160}
\definecolor{ansred}{RGB}{180,20,20}
\newcommand{\solution}{\par\noindent\textcolor{ansblue}{\textbf{解：}}}
\newcommand{\answer}{\par\noindent\textcolor{ansred}{\textbf{答案：}}}
\newcommand{\supp}[1]{\textcolor{ansred}{\textsuperscript{[补] #1}}}

\title{CS 2089M 数据库系统\\2022--2023 春夏学期 期末回忆卷}
\author{整合自两份回忆版本（CC98 论坛）\thanks{版本一：较完整的 PDF 回忆版本；版本二：选择题+判断题+大题框架回忆版本。两版为同一套试卷，此处合并整理。}}
\date{}

\begin{document}

\maketitle

\vspace{-8pt}
\begin{center}
\textbf{题型}：选择题 $8\times 3' = 24'$，判断题 $6\times 2' = 12'$，
SQL $16'$，E-R 设计 $16'$，并发控制 $8'$，ARIES $12'$，Buffer Tree $12'$。共 $100'$。
\end{center}

\section*{一、选择题（每题 3 分，共 8 题）}

\begin{enumerate}[leftmargin=*,itemsep=8pt]

\item 对于如下定义的表 (\supp{题目原文为``类似去年第二题''，以下为回忆版})：
{\small
\begin{verbatim}
create table course (
  course_id  char(2) primary key,
  credit     int,
  prereq_id  char(2),
  foreign key prereq_id references course(course_id)
    on delete cascade
    on update cascade
);
\end{verbatim}}

初始数据（\supp{来自 PDF 回忆版}）：

\begin{tabular}{c|c|c}
\texttt{course\_id} & \texttt{credit} & \texttt{prereq\_id} \\
\midrule
11 & 2 & NULL \\
22 & 4 & NULL \\
33 & 3 & NULL \\
44 & 3 & 22   \\
55 & 4 & 22   \\
66 & 3 & 33   \\
\end{tabular}

依次执行以下 SQL 语句：
{\small
\begin{enumerate}[label=(\arabic*),nosep]
  \item \texttt{INSERT INTO course VALUES ('77', 4, '11');}
  \item \texttt{DELETE FROM course WHERE course\_id = '11';}
  \item \texttt{UPDATE course SET course\_id = '11', credit = 2
          WHERE course\_id = '33';}
  \item \texttt{SELECT COUNT(*) FROM course WHERE credit > 2;}
\end{enumerate}}

最后一条语句返回的结果是：

\begin{enumerate}[label=\Alph*.]
  \item 1
  \item 2
  \item 3
  \item 4
\end{enumerate}

\solution
逐步追踪（注意 CASCADE 约束）：

\textbf{初始}：
11(2,NULL), 22(4,NULL), 33(3,NULL), 44(3,22), 55(4,22), 66(3,33)。

\textbf{(1) INSERT ('77',4,'11')}：
新增 77(4,11)。

\textbf{(2) DELETE course\_id='11'}：
删除 11。CASCADE：所有 prereq\_id=11 的记录也被删除 $\to$ 77 被删。

剩余：22(4,NULL), 33(3,NULL), 44(3,22), 55(4,22), 66(3,33)。

\textbf{(3) UPDATE course\_id='11', credit=2 WHERE course\_id='33'}：
33 改为 11(credit=2)。CASCADE：66 的 prereq\_id 跟着从 33 变为 11。

剩余：11(2,NULL)\textsuperscript{原33}, 22(4,NULL), 44(3,22), 55(4,22), 66(3,11)。

\textbf{(4) SELECT COUNT(*) WHERE credit > 2}：
credit>2 的记录：22(4), 44(3), 55(4), 66(3)，共 4 条。
11 的 credit=2 不满足。

\answer D. 4。注意 CASCADE 的连锁效应是关键。

\item 对于上题定义的表 \texttt{course}，有如下 SQL 查询：
\begin{verbatim}
SELECT c1.course_id, c1.credit, c2.course_id, c2.credit
FROM course AS c1, course AS c2
WHERE c1.prereq_id = c2.course_id
  AND c1.credit = 4
  AND c2.credit > 2;
\end{verbatim}

下列关系代数表达式中与上面的 SQL 查询\textbf{不等价}的是：
\supp{选项回忆不全，考点是：两表连接属性不同名，不能直接用自然连接}
\\[2pt]
\textcolor{gray}{⚠ 原卷选项缺失，无法确定具体选哪个。考点归纳见答案。}

\answer
核心考点：\texttt{c1.prereq\_id = c2.course\_id} 是两个\textbf{不同名}属性的等值连接，因此：
\begin{itemize}[nosep]
  \item $\pi(\sigma(\rho_{c1}(course)\times\rho_{c2}(course)))$ 形式的表达式是正确的
  \item 直接写 $\rho_{c1}(course) \Join \rho_{c2}(course)$（自然连接）\textbf{不等价}——因为 \texttt{c1} 和 \texttt{c2} 虽出自同一表，但经重命名后，自然连接会基于\textbf{所有同名属性}做 join，而题目连接条件仅涉及 \texttt{prereq\_id = course\_id} 一对不同名属性
  \item 必须显式写 $\Join_{c1.prereq\_id = c2.course\_id}$，不可用自然连接替代
\end{itemize}

\item 关系模式 $R(A,B,C,D,E,F)$ 上有函数依赖集
$F = \{A \to B,\; B \to C,\; C \to B,\; D \to E,\; D \to F\}$。
下列选项中是该关系模式的候选键的是：

\begin{enumerate}[label=\Alph*.]
  \item AC
  \item AD
  \item AE
  \item AF
\end{enumerate}

\solution
$A^+ = \{A,B,C\}$，$D^+ = \{D,E,F\}$，$AD^+ = \{A,B,C,D,E,F\} = R$。
任何不含 $D$ 的集合无法推出 $E,F$；任何不含 $A$ 的集合无法推出存在 $A,B,C$ 依赖性但已有 $B,C$ 自循环。
验证 $AD$ 是最小——$\{A,D\}$ 缺一不可。

\answer B. AD。

\item 函数依赖集 $\{A \to B,\; A \to D,\; BC \to E,\; AC \to D\}$ 的正则覆盖是：

\supp{选项回忆不全，以下为推算}

\begin{enumerate}[label=\Alph*.]
  \item $\{A \to BDE,\; BC \to E\}$
  \item $\{A \to BD,\; ABC \to E\}$
  \item $\{A \to BD,\; C \to E\}$
  \item $\{A \to BD,\; BC \to E\}$
\end{enumerate}

\solution
求正则覆盖：

\textbf{步骤 1}（RHS 已单属性）：$\{A \to B, A \to D, BC \to E, AC \to D\}$。

\textbf{步骤 2}（去掉 LHS 多余属性）：
\begin{itemize}[nosep]
  \item $AC \to D$：计算 $A^+_F = \{A,B,D\}$。$D \in A^+_F$，故 $C$ 在 $AC \to D$ 中是多余的，变为 $A \to D$（已存在）。
  \item $BC \to E$：$B^+_F = \{B\}$（无 $E$），$C^+_F = \{C\}$（无 $E$），均不多余。
\end{itemize}

\textbf{步骤 3}（去掉冗余 FD）：
得 $\{A \to B, A \to D, BC \to E\}$，即 $\{A \to BD, BC \to E\}$。

\answer D。考场上查 $BC \to E$ 中 $B$ 或 $C$ 是否多余、$C \to E$ 是否可独立（不能），可用排除法。

\item 关系 $A$ 上有 20,000 条记录。对关系 $A$ 的非键属性建立 B$^+$ 树索引。每个 B$^+$ 树节点大小为 2\,KB，每个键值大小为 32\,B，每个指针大小为 8\,B。若某次查询得到的结果包含在 5 个块中，则此次查询的开销约为：

\supp{版本二备注：``基于此建立 B+ 树非码上的索引''}

\begin{enumerate}[label=\Alph*.]
  \item 4 seek + 8 transfer
  \item 4 seek + 9 transfer
  \item 5 seek + 8 transfer
  \item 5 seek + 9 transfer
\end{enumerate}

\solution
B$^+$ 树节点容量：
\begin{itemize}[nosep]
  \item 非叶节点：$(n-1)\times 32 + n\times 8 \le 2048 \Rightarrow n \le 52$，扇出 $= 52$
  \item 叶节点：$n\times(32+8) + 8 \le 2048 \Rightarrow n \le 51$，每叶 51 条
\end{itemize}

$20,000 / 51 \approx 393$ 叶节点，$\lceil 393/52 \rceil = 8$ 个内部节点，$\to 1$ 个根。
树高 3（根 + L1 + 叶）。

查询开销（\supp{设根节点缓存在内存}）：
\begin{itemize}[nosep]
  \item 索引遍历：L1 内节点 seek 1 + transfer 1，叶节点 seek 1 + transfer 1 $\to$ 2 seek + 2 transfer
  \item 结果 5 个数据块（聚簇假设）：1 seek（定位首个块）+ 5 transfer
  \item 合计：$2+1 = 3$ seek，$2+5 = 7$ transfer。与选项均不完美匹配。
\end{itemize}

\supp{实际考试可能按另一种假设（根不缓存或结果块在索引叶中即为数据块）。因原卷数据完整但答案取决于假设，考场建议现场推。最接近的估计为 B 或 A。}

\answer 推荐按 fanout 52、高 3、根 cached、结果 5 块索引+数据合计 8 或 9 transfer 估算，选 A 或 B 均可（需视考试当场假设）。重点掌握计算过程。

\item 关系 $r$ 上有 50,000 条记录，每个块可容纳 $r$ 上的 200 条记录；关系 $s$ 上有 20,000 条记录，每个块可容纳 $s$ 上的 50 条记录。在不考虑递归分块的情况下，对 $r$ 和 $s$ 进行 hash join 的开销约为：

\begin{enumerate}[label=\Alph*.]
  \item 1950 transfer + 1350 seek
  \item 1950 transfer + 650 seek
  \item 3300 transfer + 1350 seek
  \item 3300 transfer + 650 seek
\end{enumerate}

\solution
$b_r = 50,000/200 = 250$ 块，$b_s = 20,000/50 = 400$ 块。

Hash join（不分块）：
\begin{itemize}[nosep]
  \item 分区 $r$（读+写）：$2 \times 250 = 500$ transfer
  \item 分区 $s$（读+写）：$2 \times 400 = 800$ transfer
  \item Join 阶段（读分区结果）：$250 + 400 = 650$ transfer
  \item 总 transfer：$500 + 800 + 650 = \mathbf{1950}$
\end{itemize}

Seek 次数取决于缓冲区大小。若每块一个 seek（最坏）：
$(250+400) \times 2 = 1300$ seek（写分区时每 block 一次 seek）。
加少量 join 阶段 seek $\approx$ 1350。

\answer C. 1950 transfer + 1350 seek。\supp{若假设较大的缓冲区则 seek 可降至 650 左右，选 B。实际考场需根据平时习题的缓冲区假设判断。}

\item 关系 $r$ 上有 30,000 条记录，每个块上可容纳 200 条记录。若在 $r$ 的候选键上进行等值查找，则返回结果的大小最接近：

\begin{enumerate}[label=\Alph*.]
  \item 1
  \item 20
  \item 1,500
  \item 30,000
\end{enumerate}

\solution
候选键保证唯一性——等值查找最多命中 1 条记录。

\answer A. 1。

\item 在下面的锁请求序列中：

\begin{center}
\begin{tabular}{c|c|c}
T1 & T2 & T3 \\
\midrule
lock-S(A) & & \\
& lock-X(A) & \\
& & lock-X(B) \\
& lock-X(B) & \\
& & lock-S(C) \\
& & lock-X(C) \\
\end{tabular}
\end{center}

\supp{原卷注明锁序列记不清了，写了类似的}

下列说法正确的是：

\begin{enumerate}[label=\Alph*.]
  \item 计划中没有死锁
  \item 计划中有死锁，T1 在等待 T2 放锁的同时 T2 也在等待 T1 放锁
  \item 计划中有死锁，T2 在等待 T3 放锁的同时 T3 也在等待 T2 放锁
  \item 计划中有死锁，T2 在等待 T1 放锁，T3 在等待 T2 放锁，T1 在等待 T3 放锁
\end{enumerate}

\solution
分析 wait-for 图：
\begin{itemize}[nosep]
  \item T1 持有 S(A)。T2 请求 X(A) $\to$ T2 等待 T1。
  \item T3 请求 X(B)，B 当前未被锁 $\to$ T3 获得 X(B)。
  \item T2 请求 X(B)，T3 持有 X(B) $\to$ T2 等待 T3。
  \item T3 持有 S(C)，请求 X(C) $\to$ 锁升级冲突（T3 等待自身或需先释放 S(C)），不构成跨事务等待。
\end{itemize}

Wait-for 边：T2 $\to$ T1（在 A 上），T2 $\to$ T3（在 B 上）。无环 $\to$ 无死锁。

\supp{注意：原卷锁序列为近似回忆，实际排布可能有所不同。解题关键是画 wait-for graph 判断有无环。}

\answer 在此回忆序列下选 A（无死锁）。实际考场按当时给出的序列画 wait-for 图即可。

\end{enumerate}

\newpage
\section*{二、判断题（每题 2 分，共 6 题）}

\begin{enumerate}[leftmargin=*,itemsep=6pt]

\item 任意关系模式均可分解为三范式（3NF），且分解是保持依赖的。

\solution
正确。任何关系模式都可以无损地分解为 3NF，且保持函数依赖。这是 3NF 的经典定理——与 BCNF 不同（BCNF 不一定能保持依赖）。3NF 在「消除冗余」和「保持依赖」之间做了折中。

\answer \textcolor{ansred}{$\surd$（正确）}

\item 在事务频繁更新数据的场景下，按列存储的效率比按行存储高。

\solution
错误。列存储的优势在读密集型分析查询（只需少数列时节省 I/O），对频繁更新、插入的场景反而效率低——修改一行需要更新分散在不同列存储块中的多个位置，I/O 开销远大于行存储。

\answer \textcolor{ansred}{$\times$（错误）}

\item 建立辅助索引的时候，为了降低空间开销，可以对每一个数据块建立稀疏索引。

\solution
错误。稀疏索引要求数据按搜索键\textbf{顺序存储}，通常只适用于主索引（聚簇索引）。辅助索引（非聚簇索引）的搜索键顺序与数据物理顺序不同，必须对每条记录建立索引项（稠密索引），不能使用稀疏索引。

\answer \textcolor{ansred}{$\times$（错误）}

\item 连接前尽可能早地进行投影操作可以降低投影操作的开销，所以投影操作应该尽可能早做。

\solution
正确。投影尽早做（在查询树中下推）可缩减中间结果中的列数和行数，降低后续操作（连接、排序等）的 I/O 和计算开销。这是查询优化中 ``尽早投影'' 的启发式规则。

\answer \textcolor{ansred}{$\surd$（正确）}

\item 按索引查找的时间复杂度总是比顺序查找低。\supp{版本二忘掉序号 5，此为 PDF 版题}

\solution
错误。``总是'' 过于绝对。在某些情况下（如小表、高选择性查询返回大量记录、或索引聚簇因子极差），顺序扫描可能比索引查找更快——索引查找需要先遍历 B$^+$ 树再回表读数据块，多次随机 I/O 累计可能超过一次全表顺序扫描。

\answer \textcolor{ansred}{$\times$（错误）}

\item 不能通过 2PL 产生的调度，可能可以由树协议（Tree Protocol）产生。

\solution
正确。树协议允许在访问数据项的树状结构中按特定规则加解锁（如必须先锁父节点才能锁子节点），可生成一些 2PL 无法产生的可串行化调度。两者产生的是不同的可串行化调度子集，互不包含。

\answer \textcolor{ansred}{$\surd$（正确）}

\end{enumerate}

\newpage
\section*{三、SQL 语句（16 分）}

对于如下定义的关系模式：

\begin{verbatim}
user(userID, name, gender, age, group_name)
followship(userID, followerID)
\end{verbatim}

写出可以进行如下查询的 SQL 语句：

\begin{enumerate}[label=(\arabic*),leftmargin=*]
  \item 查询所有含有男性成员的群组（group）
  \item 查询 userID 为 '1001' 的用户的追随者的信息
  \item 查出女性成员最多的群组
  \item 查询名称为 'game' 的群组当中，追随者（follower）数量高于该群组中所有用户的追随者数量的平均值的用户。
\end{enumerate}

\solution

\begin{enumerate}[label=(\arabic*),leftmargin=*]

  \item \textbf{含有男性成员的群组}：
\begin{verbatim}
SELECT DISTINCT group_name
FROM user
WHERE gender = 'male';
\end{verbatim}

  \item \textbf{1001 号用户的追随者信息}：
\begin{verbatim}
SELECT u.*
FROM user u
JOIN followship f ON u.userID = f.followerID
WHERE f.userID = '1001';
\end{verbatim}

  \item \textbf{女性成员最多的群组}：
\begin{verbatim}
SELECT group_name
FROM user
WHERE gender = 'female'
GROUP BY group_name
ORDER BY COUNT(*) DESC
LIMIT 1;
\end{verbatim}
  \supp{若有多组并列第一，需用子查询：}\\
\begin{verbatim}
SELECT group_name
FROM user
WHERE gender = 'female'
GROUP BY group_name
HAVING COUNT(*) = (
  SELECT MAX(cnt) FROM (
    SELECT COUNT(*) AS cnt FROM user
    WHERE gender = 'female' GROUP BY group_name
  ) t
);
\end{verbatim}

  \item \textbf{game 群组中追随者数高于组内均值的用户}：
\begin{verbatim}
SELECT u.userID, u.name,
       COUNT(f.followerID) AS follower_count
FROM user u
JOIN followship f ON u.userID = f.userID
WHERE u.group_name = 'game'
GROUP BY u.userID
HAVING COUNT(f.followerID) > (
  SELECT AVG(follower_count)
  FROM (
    SELECT COUNT(f2.followerID) AS follower_count
    FROM user u2
    LEFT JOIN followship f2 ON u2.userID = f2.userID
    WHERE u2.group_name = 'game'
    GROUP BY u2.userID
  ) AS avg_sub
);
\end{verbatim}

\end{enumerate}

\newpage
\section*{四、数据库设计（16 分）}

请根据以下需求设计一个跟团旅行预订平台：

\begin{itemize}[nosep]
  \item 每个用户都有用户 ID、用户名、密码和联系方式
  \item 每个旅行社需提供名称、法人代表身份证号、企业银行账号、位置，平台给每个旅行
        社分配 ID
  \item 旅行社会出售旅行计划，旅行计划的内容包含标题、时长、描述
  \item 每个用户都可以预订旅行计划，订单信息包含用户信息、旅行计划信息、订单发起
        时间、旅行计划数量、总价格、支付状态
  \item 旅行结束后用户可以对旅行社进行评价
\end{itemize}

\begin{enumerate}[label=(\arabic*),leftmargin=*]
  \item 画出 E-R 图。
  \item 将 E-R 图转化为关系模式，指明主键和外键。
\end{enumerate}

\solution

\begin{enumerate}[label=(\arabic*),leftmargin=*]

  \item \textbf{E-R 图}（\supp{文字描述，考场画图}）：

  实体：
  \begin{itemize}[nosep]
    \item \textbf{User}：属性 $\{userID, name, password, contact\}$，
          userID 为键
    \item \textbf{TravelAgency}：属性 $\{agencyID, name, legal\_ID,
          bank\_account, location\}$，agencyID 为键
    \item \textbf{TravelPlan}：属性 $\{planID, title, duration,
          description\}$，planID 为键
  \end{itemize}

  联系：
  \begin{itemize}[nosep]
    \item \textbf{Sell}：TravelAgency $\xrightarrow{1}$ Sell
          $\xrightarrow{N}$ TravelPlan（一个旅行社卖多个旅行计划）
    \item \textbf{Order}：User $\xrightarrow{M}$ Order
          $\xrightarrow{N}$ TravelPlan，带属性 $\{order\_time, quantity,
          total\_price, payment\_status\}$
    \item \textbf{Rate}：User $\xrightarrow{M}$ Rate
          $\xrightarrow{N}$ TravelAgency，带属性 $\{rating, comment,
          rate\_time\}$（旅行后评价旅行社）
  \end{itemize}

  \item \textbf{关系模式}：

  \begin{itemize}[nosep]
    \item \texttt{User(\underline{userID}, name, password, contact)}
    \item \texttt{TravelAgency(\underline{agencyID}, name, legal\_ID,
           bank\_account, location)}
    \item \texttt{TravelPlan(\underline{planID}, title, duration, description,
           \textit{agencyID})}
          \hfill FK: \texttt{agencyID $\to$ TravelAgency}
    \item \texttt{OrderInfo(\underline{orderID}, \textit{userID},
           \textit{planID}, order\_time, quantity, total\_price,
           payment\_status)}
          \hfill FK: \texttt{userID $\to$ User,\; planID $\to$ TravelPlan}
    \item \texttt{Rate(\underline{userID, agencyID}, rating, comment,
           rate\_time)}
          \hfill FK: \texttt{userID $\to$ User,\; agencyID $\to$
           TravelAgency}
  \end{itemize}

\end{enumerate}

\newpage
\section*{五、并发控制（8 分）}

\supp{原卷注明序列为近似回忆，以下为回忆序列}

对于如下的访问序列：
\begin{center}
\texttt{r1(A) w2(A) w2(B) w3(A) r1(A) w3(B) r4(C) r4(B) w5(C) w5(B) r3(C)}
\end{center}

其中 $\texttt{ri(D)}$ 表示事务 $T_i$ 对数据 D 进行读操作，$\texttt{wi(D)}$ 表示写操作。

\begin{enumerate}[label=(\arabic*),leftmargin=*]
  \item 根据访问序列画出各个事务的前驱图（precedence graph）
  \item 判断这一访问序列是否冲突可串行化并说明理由
  \item 判断这一访问序列是否可以由两阶段锁协议（2PL）产生并说明理由
\end{enumerate}

\solution

\begin{enumerate}[label=(\arabic*),leftmargin=*]

  \item \textbf{前驱图}：找出所有冲突操作对（不同事务访问同一数据且至少一个是写）。

  \begin{itemize}[nosep]
    \item \texttt{r1(A)} $\to$ \texttt{w2(A)}：$T_1 \to T_2$
    \item \texttt{w2(A)} $\to$ \texttt{w3(A)}：$T_2 \to T_3$
    \item \texttt{w2(A)} $\to$ \texttt{r1(A)}（第二条）：$T_2 \to T_1$
    \item \texttt{w3(A)} $\to$ \texttt{r1(A)}（第二条）：$T_3 \to T_1$
    \item \texttt{w2(B)} $\to$ \texttt{w3(B)}：$T_2 \to T_3$
    \item \texttt{w2(B)} $\to$ \texttt{r4(B)}：$T_2 \to T_4$
    \item \texttt{w3(B)} $\to$ \texttt{r4(B)}：$T_3 \to T_4$
    \item \texttt{r4(C)} $\to$ \texttt{w5(C)}：$T_4 \to T_5$
    \item \texttt{r4(B)} $\to$ \texttt{w5(B)}：$T_4 \to T_5$
    \item \texttt{w5(B)} $\to$ （无关）\texttt{r3(C)}；
          \texttt{w5(C)} $\to$ \texttt{r3(C)}：$T_5 \to T_3$
  \end{itemize}

  化简边：$T_1 \to T_2 \to T_3 \to T_1$（环），$T_1 \to T_2 \to T_1$（直接环）。

  \supp{完整前驱图：T1-T2-T3 三者两两成环，T4、T5 为下游。}

  \item \textbf{不可串行化}。前驱图存在环（$T_1 \to T_2 \to T_1$ 或
        $T_1 \to T_2 \to T_3 \to T_1$），不满足冲突可串行化的充要条件
        （前驱图无环）。

  \item \textbf{不可由 2PL 产生}。2PL 产生的所有调度都是冲突可串行化的
        （前驱图一定无环）。既然该序列不可串行化，就必然不能由 2PL 产生。

\end{enumerate}

\newpage
\section*{六、ARIES 恢复算法（12 分）}

\supp{原卷注明日志为近似回忆}

给出如下的 ARIES 日志：

{\small
\begin{verbatim}
LSN  TID  Type   PageID  Slot  Before  After  prevLSN
1    T1   upd    6601    1     A       B      0
2    T2   upd    2628    2     C       D      0
3    T1   upd    2628    1     X       Y      1
4    T3   upd    6601    2     E       F      0
5    T1   upd    6601    3     G       H      3
6    T2   CLR                                ← checkpoint
7    T1   upd    ...    ...   ...    ...    5
8    T2   upd    ...    ...   ...    ...    6
9    T2   commit                              8
10   T3   upd    6601    2     M       N      4
\end{verbatim}}

\supp{具体 LSN、PageID、Slot、值均无法完整回忆，以上日志结构仅用于说明考点}

\begin{enumerate}[label=(\arabic*),leftmargin=*]
  \item 写出在分析过程（Analysis Phase）中被加入到脏页表（Dirty Page Table）的记录
  \item 重做过程从哪一条日志开始？
  \item 写出恢复过程中增加的日志（要求包含事务号、修改页号、slot 号和修改的值）
  \item 恢复过程结束后，6601.1 和 2628.2 的值分别是多少？
\end{enumerate}

\solution
\supp{以下解答基于 ARIES 标准流程，因原卷具体日志数据缺失，无法给出确切数值。}

\begin{enumerate}[label=(\arabic*),leftmargin=*]

  \item \textbf{脏页表}：Analysis 阶段从 checkpoint 开始向前扫描，将 checkpoint 时活跃事务涉及的所有已修改页加入脏页表，再加上 checkpoint 之后的所有 update log 涉及的页。
  此处 checkpoint 前涉及页：6601, 2628；checkpoint 后 LSN 7--10 也涉及 6601，均入表。

  \item \textbf{Redo 起点}：脏页表中所有页的最早修改 LSN。
  6601 最早修改为 LSN 1，2628 最早修改为 LSN 2，取最小 = \textbf{LSN 1}。

  \item \textbf{新增 CLR 日志}：Undo 阶段对 loser（未 commit 的 T3）的每条 update，
  按 prevLSN 链反向 undo。对每条 undo 写一条 CLR（Compensation Log Record），
  含 undo 后的 before 值。\supp{具体值依赖缺失的日志数据。}

  \item \textbf{最终值}：Redo 阶段从 LSN 1 重放到末尾（包括 abort 产生的 CLR），
  所有已 commit 事务（T1, T2）的更新生效，T3 被 undo 撤销。
  6601.1 最终为 Redo 后值，2628.2 同理。\supp{具体值依赖缺失的日志数据。}

  \answer 考点：Analysis 确定脏页表+loser/winner 集合；Redo 从最早脏页 LSN 起；
  Undo 沿 prevLSN 链反向撤销 loser；CLR 格式为补偿日志。答案形式对应实际日志。

\end{enumerate}

\newpage
\section*{七、Buffer Tree（12 分）}

Buffer Tree 是一种写优化的索引结构——在 B$^+$ 树的每个内节点处增加一块缓冲区。
插入键值时，先写入根节点缓冲区。当某节点缓冲区溢出时：
\begin{itemize}[nosep]
  \item 若子节点为内节点：将缓冲区所有元素下放到子节点缓冲区中
  \item 若子节点为叶子：将缓冲区所有元素插入子节点，子节点溢出时正常分裂
\end{itemize}

本题讨论缓冲区大小为 2 且度 $n=4$ 的 Buffer Tree。初始状态如题图。

\begin{enumerate}[label=(\arabic*),leftmargin=*]
  \item 现对该树做范围查询（不小于 30 且不大于 50），共有多少节点被访问？
  \item 向该树中依次插入三个键值：38、92、28。画出三次插入操作之后的 Buffer Tree。
  \item 在第 (2) 问的插入操作中，每次插入分别有多少节点被修改（分裂新增节点视同修改；
        根节点始终缓存于内存中，不计入）？
\end{enumerate}

\solution

\supp{原卷附有初始树图。因图片无法重现，以下基于 B$^+$ 树 $n=4$（每节点最多 4 个指针，3 个键）、buffer size=2 的标准 Buffer Tree 行为进行推演。}

\begin{enumerate}[label=(\arabic*),leftmargin=*]

  \item \textbf{范围查询 [30, 50]}：从根沿树下降到叶子，遍历叶子链表。
  设初始树高为 $h$（含根节点），需访问根 + $h-2$ 个内部节点 +
  在范围 [30, 50] 内的所有叶子节点。\supp{具体节点数取决于初始树图。}

  \item \textbf{插入 38, 92, 28}：

  \begin{itemize}[nosep]
    \item 插入 38：38 写入根 buffer（buffer 未满，不溢出）\\
    \item 插入 92：92 写入根 buffer $\to$ 根 buffer 满（2 个元素），触发下放。
    92 被下放到对应子节点（假设右子树）buffer，沿途 buffer 可能下放至叶子。
    若叶子满则分裂。
    \item 插入 28：28 写入根 buffer，触发下放到左子树，可能继续下放或分裂。
  \end{itemize}

  最终形态：插入过程可能产生叶子分裂（因 $n=4$，叶满 4 元素时分裂为两个 2-元素的节点），
  并有部分键值滞留于各级 buffer 中。

  \supp{因原卷有图，实际画图需根据初始 B$^+$ 树的具体键值和结构。此处仅描述推演过程。}

  \item \textbf{每次插入修改的节点数（不计根）}：

  \begin{itemize}[nosep]
    \item 插入 38：仅写根 buffer，不触及磁盘节点 $\to$ 修改 0 个节点
    \item 插入 92：根 buffer 满 $\to$ 下放，修改被下放经过的各层节点及叶子，
          若叶子分裂则新增节点亦计数 $\to$ \supp{取决于树结构}
    \item 插入 28：同理，触发下放和可能的分裂 $\to$ \supp{取决于树结构}
  \end{itemize}

  \answer 考点：Buffer Tree 通过集中批量下放减少磁盘 I/O，insert 不计根。
  解题关键是追踪 buffer 溢出触发下放的链式反应，以及叶子分裂时的新节点分配。

\end{enumerate}

\vfill
\hrule
\vspace{4pt}
{\footnotesize \textcolor{ansred}{⚠ 本卷由两份 CC98 论坛回忆版本整合而成。选择题第 1--2、4--8 题来自 PDF 版（较完整）；第 3 题和判断题来自两版互校；大题综合两版。标注 \textsuperscript{[补]} 的内容为整理者根据课内知识点对缺失数据的技术性补充，非原卷内容。}}

\end{document}
